\documentclass[aspectratio=169,11pt]{beamer}
\usetheme{default}
\setbeamertemplate{navigation symbols}{}
\usepackage[T1]{fontenc}
\usepackage{lmodern}
\usepackage{booktabs}
\usepackage{siunitx}
\usepackage{multirow}
\usepackage{graphicx}
\usepackage{xcolor}
\usepackage{amsmath}
\usepackage{tikz}

% ---- data-series colours (defined BEFORE the theme, which references catk) ----
\definecolor{cfixed}{RGB}{54,100,158}
\definecolor{cvt}{RGB}{206,123,26}
\definecolor{catk}{RGB}{178,45,45}
\newcommand{\FIX}[1]{\textcolor{cfixed}{#1}}
\newcommand{\VT}[1]{\textcolor{cvt}{#1}}
\newcommand{\ATK}[1]{\textcolor{catk}{#1}}

% ===================== GREEN THEME =====================
\definecolor{thHeader}{RGB}{44,70,40}
\definecolor{thMid}{RGB}{89,127,60}
\definecolor{thTint}{RGB}{232,239,223}
\definecolor{thSubtitle}{RGB}{205,222,193}
\definecolor{thGray}{RGB}{120,120,120}

\setbeamercolor{normal text}{fg=black,bg=white}
\setbeamercolor{frametitle}{bg=thHeader,fg=white}
\setbeamercolor{framesubtitle}{fg=thSubtitle}
\setbeamercolor{title}{fg=thHeader}
\setbeamercolor{subtitle}{fg=thMid}
\setbeamercolor{author}{fg=thGray}
\setbeamercolor{date}{fg=thGray}
\setbeamercolor{institute}{fg=thGray}
\setbeamercolor{section in head/foot}{fg=thGray,bg=white}
\setbeamercolor{block title}{bg=thMid,fg=white}
\setbeamercolor{block body}{bg=thTint,fg=black}
\setbeamercolor{block title alerted}{bg=catk,fg=white}
\setbeamercolor{block body alerted}{bg=catk!8,fg=black}
\setbeamercolor{itemize item}{fg=thMid}
\setbeamercolor{itemize subitem}{fg=thMid}
\setbeamercolor{itemize subsubitem}{fg=thMid}
\setbeamercolor{item projected}{bg=thMid,fg=white}
\setbeamercolor{footline}{fg=thGray,bg=white}
\setbeamercolor{page number in head/foot}{fg=thGray}
\setbeamercolor{section number projected}{bg=thMid,fg=white}
\setbeamercolor{caption}{fg=thGray}

\setbeamerfont{frametitle}{size=\large,series=\bfseries}
\setbeamerfont{framesubtitle}{size=\scriptsize,series=\mdseries}
\setbeamerfont{section in head/foot}{size=\tiny}
\setbeamerfont{block title}{series=\bfseries}
\setbeamerfont{footline}{size=\scriptsize}

\setbeamertemplate{itemize item}{\tikz\fill[thMid] (0,0) circle (0.42ex);}
\setbeamertemplate{itemize subitem}{\tikz\fill[thMid] (0,0) circle (0.34ex);}
\setbeamertemplate{itemize subsubitem}{\tikz\draw[thMid,line width=0.6pt] (0,0) circle (0.34ex);}

\newcommand{\thNumCircle}[1]{%
  \tikz[baseline=-0.65ex]\node[circle,fill=thMid,text=white,inner sep=0pt,
       minimum size=2.1ex,font=\scriptsize\bfseries]{#1};%
}
\setbeamertemplate{enumerate item}{\thNumCircle{\insertenumlabel}}
\setbeamertemplate{enumerate subitem}{\thNumCircle{\insertsubenumlabel}}
\setbeamertemplate{enumerate subsubitem}{\thNumCircle{\insertsubsubenumlabel}}

\setbeamertemplate{section in toc}{%
  \leavevmode\thNumCircle{\inserttocsectionnumber}\hskip0.6em{\color{black}\inserttocsection}\par\vskip1.1ex}
\setbeamertemplate{subsection in toc}{%
  \leavevmode\leftskip=3.2em\tikz\fill[thMid] (0,0) circle (0.34ex);\hskip0.6em{\color{black}\inserttocsubsection}\par\vskip0.4ex}

\defbeamertemplate*{section in head/foot}{green theme}{\textcolor{thMid}{\bfseries\insertsectionhead}}
\defbeamertemplate*{section in head/foot shaded}{green theme}{\textcolor{thGray}{\insertsectionhead}}
\setbeamertemplate{headline}{%
  \begin{beamercolorbox}[wd=\paperwidth,ht=3.4ex,dp=1.1ex]{section in head/foot}
    \insertsectionnavigationhorizontal{\paperwidth}{\hskip0.4cm}{\hskip0.4cm}
  \end{beamercolorbox}%
  {\color{thHeader}\rule{\paperwidth}{1.1pt}}%
}

\setbeamertemplate{frametitle}{%
  \nointerlineskip
  \begin{beamercolorbox}[wd=\paperwidth,leftskip=0.5cm,rightskip=0.5cm,vmode]{frametitle}
    \vspace{1.4ex}
    \usebeamerfont{frametitle}\insertframetitle\par
    {\usebeamerfont{framesubtitle}\usebeamercolor[fg]{framesubtitle}\insertframesubtitle\strut\par}
    \vspace{0.9ex}
  \end{beamercolorbox}%
}

\setbeamertemplate{footline}{%
  {\color{thMid}\hrule width\paperwidth height0.8pt}%
  \begin{beamercolorbox}[wd=\paperwidth,ht=2.4ex,dp=1.1ex,leftskip=0.5cm,rightskip=0.5cm]{footline}
    \usebeamerfont{footline}\insertshorttitle\hfill\insertframenumber\,/\,\inserttotalframenumber
  \end{beamercolorbox}%
}

\setbeamertemplate{title page}{%
  \vbox to \textheight{%
    \vspace*{0.5cm}
    {\color{thMid}\rule{\linewidth}{2.4pt}}\\[1.4cm]
    {\usebeamercolor[fg]{title}\usebeamerfont{title}\inserttitle\par}
    \vspace{0.5cm}
    {\usebeamercolor[fg]{subtitle}\usebeamerfont{subtitle}\insertsubtitle\par}
    \vfill
    {\usebeamercolor[fg]{author}\usebeamerfont{author}\insertauthor\par}
    \vspace{0.2cm}
    {\usebeamercolor[fg]{date}\usebeamerfont{date}\insertdate\par}
    \vspace{0.6cm}
    {\color{thMid}\rule{\linewidth}{2.4pt}}
    \vfill
  }
}
\setbeamerfont{title}{size=\Large,series=\bfseries}
\setbeamerfont{subtitle}{size=\normalsize,series=\mdseries}
% =================== END GREEN THEME ====================

\newcommand{\figorbox}[2]{%
  \IfFileExists{#1}{\includegraphics[width=#2]{#1}}%
  {\fbox{\parbox[c][3.0cm][c]{#2}{\centering\small\ATK{[FIGURE PLACEHOLDER]}\\\texttt{\detokenize{#1}}}}}}
\graphicspath{{figs/}{./}}
\sisetup{detect-all, table-number-alignment=center}
\title[VT-SNN Defense]{Variable Timing in Spiking Neural Networks:\\ Efficiency Gains and the Attack Surface They Create}
\subtitle{BSc Thesis Defense}
\author{Rohan}
\date{\today}
\begin{document}
\begin{frame}\titlepage\end{frame}
\begin{frame}{Outline}\framesubtitle{Roadmap}\tableofcontents\end{frame}

\section{Background}
\begin{frame}{Motivation}\framesubtitle{Why adaptive inference in SNNs?}
\begin{itemize}
  \item SNNs process input over $T$ discrete timesteps; cost scales with $T$.
  \item Most samples are easy, and are classified correctly long before $T=8$.
  \item \VT{Variable Timing (VT)}: halt once confidence crosses a threshold $\theta$.
  \item Promise: large compute savings, with no change to the backbone.
\end{itemize}
\vspace{0.4em}
\begin{block}{The question this thesis asks}
VT is almost always compared against \FIX{full-horizon $T=8$} inference. That is not the
fair baseline. What happens under a \FIX{fixed-$T$ control at the same compute budget}?
\end{block}
\end{frame}

\begin{frame}{Spiking Neural Networks}\framesubtitle{Primer}
\begin{columns}[T]
\begin{column}{0.55\textwidth}
\begin{itemize}
  \item Leaky integrate-and-fire neurons, direct input encoding.
  \item Binary spikes turn multiplications into accumulations.
  \item Cost metric: \textbf{SynOps} = dense ops $\times$ activation sparsity.
  \item Deterministic at inference: no membrane noise.
\end{itemize}
\end{column}
\begin{column}{0.40\textwidth}
\centering\figorbox{SNN.png}{\linewidth}
\end{column}
\end{columns}
\end{frame}

\begin{frame}{The VT Mechanism}\framesubtitle{}
\centering\figorbox{VT.png}{0.70\linewidth}
\end{frame}

\begin{frame}{Related Work}\framesubtitle{Three threads, studied separately}
\small
\begin{columns}[T]
\begin{column}{0.5\textwidth}
\begin{block}{Adaptive Inference}
BranchyNet, MSDNet, ACT; Dynamic Confidence, DT-SNN, SpikeCP, VT-SNN.
\end{block}
\vspace{0.4em}
\begin{block}{SNN Efficiency}
SynOps accounting, timestep reduction, AutoSNN.
\end{block}
\end{column}
\begin{column}{0.5\textwidth}
\begin{block}{SNN Robustness}
Inherent-robustness claims, certified training, neuromorphic temporal coding.
\end{block}
\vspace{0.4em}
\begin{block}{Reliable Evaluation}
PGD alone is insufficient; AutoAttack; adaptive attacks.
\end{block}
\end{column}
\end{columns}
\end{frame}

\begin{frame}{The Gap}\framesubtitle{What prior work controls for}
\centering\small
\begin{tabular}{lcccc}
\toprule
 & Compute-matched & Adaptive & Energy as & Calib./eval \\
 & baseline & attacks & attack surface & separation \\
\midrule
DT-SNN / Dynamic Confidence & \ATK{no} & \ATK{no} & \ATK{no} & \ATK{no} \\
Neuromorphic robustness     & \ATK{no} & partial  & \ATK{no} & \ATK{no} \\
SNN energy accounting       & \ATK{no} & \ATK{no} & \ATK{no} & n/a \\
SpikeCP / VT-SNN            & \ATK{no} & \ATK{no} & \ATK{no} & partial \\
\midrule
\textbf{This thesis}        & \textbf{yes} & \textbf{yes} & \textbf{yes} & \textbf{yes} \\
\bottomrule
\end{tabular}
\vspace{0.6em}
\begin{block}{}
\small No prior VT study compares against a \FIX{compute-matched} fixed-$T$ control.
\end{block}
\end{frame}

\begin{frame}{Problem Statement}
\framesubtitle{}
\begin{itemize}
    \item \textbf{VT is reported as a major efficiency improvement} for
    Spiking Neural Networks (SNNs).
    \item However, existing comparisons often use
    \textbf{full-horizon inference} as the baseline, rather than a
    \textbf{fixed-time model at the same compute budget}.
    \item Therefore, it remains unclear whether VT provides a
    \textbf{genuine efficiency advantage} or simply benefits from a
    favorable comparison.
    \item Moreover, VT introduces a \textbf{data-dependent halting decision},
    creating a new control channel whose vulnerability to
    \textbf{adversarial manipulation} has not been sufficiently studied.
\end{itemize}
\vspace{0.8mm}
\begin{alertblock}{This Thesis}
Evaluates VT from both \textbf{compute-matched efficiency} and
\textbf{adversarial robustness} perspectives.
\end{alertblock}
\end{frame}

\begin{frame}{Contributions}
\framesubtitle{Four Key Contributions}
\begin{columns}[T,onlytextwidth]
\begin{column}{0.48\textwidth}
\begin{block}{1. Compute-Matched Evaluation}
\small
A \textbf{compute-matched evaluation protocol} for VT,
including calibration/test separation.
\end{block}
\vspace{8mm}
\begin{block}{3. Selection-Bias Analysis}
\small
Identification of \textbf{selection bias} in the standard
VT operating-point criterion, with a neutral replacement.
\end{block}
\end{column}
\begin{column}{0.48\textwidth}
\begin{block}{2. VT Advantage}
\small
Evidence that the VT advantage is \textbf{real but small},
and \textbf{contingent on backbone training}.
\end{block}
\vspace{4mm}
\begin{block}{4. Compute-Overuse Attack}
\small
A \textbf{compute-overuse attack} demonstrating that the
halting gate is an \textbf{energy attack surface} across
three backbones.
\end{block}
\end{column}
\end{columns}
\end{frame}

\begin{frame}{Research Questions}\framesubtitle{}
\begin{block}{RQ1 --- Efficiency}
Does VT outperform a \FIX{compute-matched fixed-$T$} baseline on clean accuracy, and
does that hold across architectures and training regimes?
\end{block}
\begin{block}{RQ2 --- Robustness}
Does any efficiency advantage survive adversarial evaluation, and does the halting
mechanism itself constitute an exploitable attack surface?
\end{block}
\end{frame}

\section{Method}
\begin{frame}{Compute-Matched Framework}\framesubtitle{The core methodological move}
\begin{enumerate}
  \item Sweep \FIX{fixed-$T$} for $T=1 \ldots 8$; record accuracy and SynOps.
  \item Sweep \VT{VT} over a threshold grid; record accuracy, stopping depth, SynOps.
  \item For each VT point, find the fixed-$T$ control with the \textbf{nearest SynOps}.
  \item Report the difference, not VT versus $T=8$.
\end{enumerate}
\begin{alertblock}{Why this matters}
Comparing VT against $T=8$ guarantees a win. Equal-compute comparison is the only test
of the adaptive-inference hypothesis.
\end{alertblock}
\end{frame}

\begin{frame}{Operating-Point Selection}\framesubtitle{A bias we found, and fixed}
\begin{columns}[T]
\begin{column}{0.5\textwidth}
\textbf{\ATK{Biased (standard)}}\\
\vspace{3mm}
Pick $\theta$ maximising (VT $-$ matched fixed).
\begin{itemize}\small
  \item Optimises the reported quantity.
  \item Seeks wherever the baseline is weakest.
  \item Selected a point \textbf{31 pts below} full horizon in one configuration.
\end{itemize}
\end{column}
\begin{column}{0.5\textwidth}
\textbf{Neutral (ours)}\\
\vspace{3mm}
Lowest-SynOps $\theta$ whose calibration accuracy is within $\delta$ of full horizon.
\begin{itemize}\small
  \item Cannot flatter VT by construction.
  \item Standard early-exit convention.
  \item Stable for $\delta \in \{0.5, 1.0, 2.0\}$.
\end{itemize}
\end{column}
\end{columns}
\vspace{0.5em}
\begin{block}{Effect on our own prior numbers}
\small ResNet-18 clean $+4.19 \rightarrow \mathbf{+2.73}$ \qquad
ResNet-18 adversarial $+2.13 \rightarrow \mathbf{+1.37}$
\end{block}
\end{frame}

\begin{frame}{Threat Models}\framesubtitle{Two distinct questions}
\begin{columns}[T]
\begin{column}{0.5\textwidth}
\textbf{Label robustness}
\vspace{3mm}\\ Does the prediction stay correct?
\begin{itemize}\small
  \item PGD-20 untargeted
  \item AutoAttack (standard ensemble)
  \item wrong-early-exit (exit-aware)
\end{itemize}
\end{column}
\begin{column}{0.5\textwidth}
\textbf{\ATK{Compute robustness}}
\vspace{3mm}\\ Does inference \emph{cost} stay bounded?
\begin{itemize}\small
  \item compute-overuse objective
  \item maximise executed steps
  \item while \emph{preserving} the label
\end{itemize}
\end{column}
\end{columns}
\vspace{0.5em}
\begin{alertblock}{}
\small A VT system can be attacked without ever being wrong. Standard robustness
evaluation cannot see this.
\end{alertblock}
\end{frame}

\begin{frame}{The Compute-Overuse Attack}\framesubtitle{Attacking the halting gate}
\begin{itemize}
  \item Objective: hold per-step confidence \emph{below} $\theta$ so the model never halts early.
  \item Simultaneously preserve the true label, so outputs still look correct.
  \item $\ell_\infty$, $\varepsilon = 8/255$, PGD-20, full temporal trace exposed to the attacker.
\end{itemize}
\begin{block}{Consequence}
Energy and latency become adversarially controllable while accuracy is untouched. The
victim sees correct predictions and a larger bill.
\end{block}
\end{frame}

\begin{frame}{Dataset and Backbones}\framesubtitle{Setup (1/3)}
\begin{block}{Dataset}
\textbf{CIFAR-10}, with train / val / calibration / test splits. Test used only for final reporting.
\end{block}
\vspace{0.3em}
\begin{block}{Backbones (deterministic, direct encoding, $T_{\max}=8$)}
\begin{itemize}\small
  \item ResNet-18, four residual stages, 11.2M parameters
  \item VGG16, five stages $(2,2,3,3,3)$, 33.6M parameters
  \item MobileNetV2, inverted residuals, 2.2M parameters
\end{itemize}
\end{block}
\vspace{0.3em}
\begin{alertblock}{Note}
\small No stochastic noise, no dynamic sigma, no Parseval regularisation.
\end{alertblock}
\end{frame}

\begin{frame}{Training Setup}\framesubtitle{Setup (2/3)}
\begin{block}{Optimization}
\small SGD, base LR $0.1$, momentum $0.9$, weight decay $10^{-4}$, 5 warmup epochs, 100 epochs.
Mixed precision, gradient clipping $1.0$, cosine schedule.
\end{block}
\vspace{0.4em}
\begin{columns}[T]
\begin{column}{0.5\textwidth}
\begin{block}{Two Training Regimes}
\small\textit{(pure ratios, no annealing)}
\begin{itemize}\small
  \item \textbf{Clean:} 100\% clean examples
  \item \textbf{Adversarial:} 100\% RFGSM adversarial examples, $\varepsilon = 8/255$
\end{itemize}
\end{block}
\end{column}
\begin{column}{0.5\textwidth}
\begin{block}{Two Timestep Schedules}
\begin{itemize}\small
  \item \textbf{Fixed-$T8$:} train only at $T=8$
  \item \textbf{Mixed-$T$:} train across $T=1 \ldots 8$
\end{itemize}
\end{block}
\end{column}
\end{columns}
\end{frame}

\begin{frame}{Attack Setup}\framesubtitle{Setup (3/3)}
\begin{block}{Attack Configuration}
\begin{itemize}
  \item $\ell_\infty$, $\varepsilon = 8/255$, PGD step $2/255$, 20 steps, random start.
  \item AutoAttack: standard ensemble (APGD-CE, APGD-T, FAB-T, Square).
  \item \textbf{Adaptive through the halting decision:} gradients propagate through the
        full temporal trace, so the attacker sees and can steer every per-step confidence.
  \item Class-balanced attack subsets; $n$ is reported with every number.
\end{itemize}
\end{block}
\vspace{0.4em}
\begin{alertblock}{Key Design Point}
\small This is the key design point for RQ2. The attack is not a label attack applied
to a VT model; it differentiates through the exit gate.
\end{alertblock}
\end{frame}

\section{Results}
\begin{frame}{Results}\framesubtitle{}\centering\Large Results\end{frame}

\begin{frame}{RQ1: The Accuracy--Compute Frontier}\framesubtitle{Representative configuration: VGG16, clean-trained}
\centering\figorbox{frontier_vgg_clean.png}{0.60\linewidth}
\\\small
\VT{VT} sits above the \FIX{fixed-$T$} curve at equal SynOps: the adaptive policy is not
merely reproducing a static truncation depth.
\end{frame}

\begin{frame}{RQ1: Frontiers Across All Configurations}\framesubtitle{Small multiples: 3 backbones $\times$ 2 regimes}
\centering\figorbox{frontier_small_multiples.png}{0.65\linewidth}
\\\small
Shared axes. The same shape recurs in all six panels, including \textbf{MobileNetV2},
whose depthwise-separable convolutions give a different SynOps profile. The effect is
not a dense-convolution artifact.
\end{frame}

\begin{frame}{RQ1: Selected Operating Points}\framesubtitle{Neutral criterion $\delta=1.0$; selected on calibration, reported on test}
\centering\small
\begin{tabular}{llS[table-format=1.2]S[table-format=1.2]S[table-format=2.2]cS[table-format=2.2]S[table-format=+1.2]}
\toprule
Backbone & Regime & {$\theta$} & {Depth} & {\VT{VT} acc.} & {Matched} & {\FIX{Fixed} acc.} & {$\Delta$} \\
\midrule
ResNet-18   & Clean & 0.90 & 3.01 & 90.62 & $T{=}3$ & 87.89 & +2.73 \\
VGG16       & Clean & 0.80 & 3.89 & 87.75 & $T{=}4$ & 84.63 & +3.12 \\
MobileNetV2 & Clean & 0.99 & 7.41 & 78.08 & $T{=}7$ & 74.71 & +3.37 \\
\midrule
ResNet-18   & Adv.  & 0.60 & 5.01 & 72.83 & $T{=}5$ & 71.46 & +1.37 \\
VGG16       & Adv.  & 0.50 & 5.74 & 64.64 & $T{=}6$ & 63.75 & +0.89 \\
MobileNetV2 & Adv.  & 0.60 & 7.21 & 50.92 & $T{=}7$ & 50.64 & +0.28 \\
\bottomrule
\end{tabular}
\vspace{0.4em}
\begin{block}{Reading}
\small \textbf{Clean-trained: a consistent $\approx +3$ across all three backbones.}
Adversarially trained: much smaller, and the lower two rows sit within the noise floor
($\pm 0.9$ to $1.0$ at these $n$). The depth column matters: MobileNetV2 saves only
about 7\% of steps, so its $+3.37$ is near-full-horizon inference.
\end{block}
\end{frame}

\begin{frame}{RQ2: Under PGD, VT Looks Fine}\framesubtitle{Adversarially trained models, PGD-20}
\centering\small
\begin{tabular}{lS[table-format=2.2]S[table-format=2.2]S[table-format=+2.2]}
\toprule
Backbone & {\FIX{Fixed $T{=}8$}} & {\VT{VT selected}} & {$\Delta$} \\
\midrule
ResNet-18   & 46.10 & 48.70 & +2.60 \\
VGG16       & 45.90 & 44.45 & -1.45 \\
MobileNetV2 & 35.60 & 34.75 & -0.85 \\
\bottomrule
\end{tabular}
\vspace{0.7em}
\begin{block}{Interim conclusion, revised on the next slide}
\small VT is competitive under PGD, and on ResNet-18 it is ahead. A study reporting only
PGD would conclude that VT is robustness-neutral.
\end{block}
\end{frame}

\begin{frame}{RQ2: Under AutoAttack, the Picture Changes}\framesubtitle{Same checkpoints, same $\varepsilon$, stronger attack}
\centering\small
\begin{tabular}{lS[table-format=2.2]S[table-format=2.2]S[table-format=+2.2]S[table-format=+2.2]}
\toprule
Backbone & {\FIX{Fixed $T{=}8$}} & {\VT{VT selected}} & {$\Delta$ (AA)} & {$\Delta$ (PGD)} \\
\midrule
ResNet-18   & 28.85 & 24.15 & -4.70 & +2.60 \\
VGG16       & 10.75 &  8.55 & -2.20 & -1.45 \\
MobileNetV2 & 27.60 & 27.55 & -0.05 & -0.85 \\
\bottomrule
\end{tabular}
\vspace{0.5em}
\begin{alertblock}{The reversal, stated precisely}
\small On \textbf{ResNet-18} the sign flips between attacks: $+2.60$ ahead under PGD,
\ATK{$-4.70$} behind under AutoAttack. \textbf{These are single runs, and the AutoAttack
result does not replicate}: over three training seeds the same comparison gives
$+2.73 \pm 2.29$ on ResNet-18 and $-2.03 \pm 4.97$ on VGG16. The apparent reversal is a
property of one checkpoint, not of VT.
\end{alertblock}
\end{frame}

\begin{frame}{RQ2: The Gap Across All Configurations}\framesubtitle{Adversarially trained; same checkpoints, two attacks}
\centering\figorbox{pgd_aa_gap.png}{0.80\linewidth}
\begin{alertblock}{}
\small Single-run result. Across three training seeds these AutoAttack differences do not
replicate: VT is ahead on ResNet-18 ($+2.73$) and behind on VGG16 ($-2.03$), with the
sign set by the training seed.
\end{alertblock}
\end{frame}

\begin{frame}{RQ2: Against the Compute-Matched Control}\framesubtitle{Sharpest form of the result: ResNet-18 adversarial, $n=1000$}
\centering\small
\begin{tabular}{lS[table-format=2.1]S[table-format=2.1]}
\toprule
Operating point & {PGD-20} & {AutoAttack} \\
\midrule
\FIX{Fixed $T=8$}                  & 44.1 & 22.2 \\
\FIX{Matched fixed $T=4$}          & 40.2 & \bfseries 33.7 \\
\VT{VT selected ($\theta{=}0.50$)} & 43.1 & \bfseries 10.5 \\
\bottomrule
\end{tabular}
\vspace{0.6em}
\begin{alertblock}{}
\small Under PGD, VT (43.1) beats its matched control (40.2). Under AutoAttack the same
comparison inverts by \textbf{23 points} (10.5 versus 33.7). Under strong attack the best
fixed-time deployment is a \emph{short} horizon, not $T=8$, and not VT.
\end{alertblock}
\end{frame}

\begin{frame}{Do the Gains Survive Training Seeds?}\framesubtitle{Three independently trained seeds per backbone; paired within each checkpoint}
\centering\small
\begin{tabular}{lcccc}
\toprule
 & \multicolumn{2}{c}{ResNet-18 adv.} & \multicolumn{2}{c}{VGG16 adv.} \\
\cmidrule(lr){2-3}\cmidrule(lr){4-5}
Condition & Mean $\pm$ std & Verdict & Mean $\pm$ std & Verdict \\
\midrule
Clean       & $+2.13 \pm 0.85$ & \textbf{real}  & $+2.80 \pm 1.18$ & \textbf{real} \\
PGD-20      & $+0.60 \pm 1.20$ & noise          & $-0.13 \pm 1.33$ & noise \\
AutoAttack  & $+2.73 \pm 2.29$ & \ATK{marginal} & $-2.03 \pm 4.97$ & noise \\
\bottomrule
\end{tabular}
\vspace{0.5em}
\begin{block}{What replicates}
\small The \textbf{clean advantage replicates on both backbones}, in both cases larger
than twice its own standard deviation. No attacked comparison does.
\end{block}
\begin{alertblock}{What does not replicate: the AutoAttack result}
\small Per-seed AutoAttack deltas are $+0.90, +5.30, +2.00$ on ResNet-18 and
$-3.50, -6.10, +3.50$ on VGG16. \textbf{The sign depends on the training seed.} A
single-run study of either backbone would confidently report an effect that is not there.
\end{alertblock}
\end{frame}

\begin{frame}{Compute-Overuse: Energy as an Attack Surface}\framesubtitle{Depth inflation at the selected operating point, label preserved}
\centering\figorbox{compute_overuse.png}{0.70\linewidth}
\begin{alertblock}{Architecture-independent}
\small Clean-trained models inflate by $\approx +1.15$ steps in all three backbones;
adversarially trained models by $<0.3$. Adversarial training is an unintentional defence
against energy attacks.
\end{alertblock}
\end{frame}

\begin{frame}{The Compute-Overuse Attack}\framesubtitle{Energy as an attack surface, all three backbones}
\centering\small
\begin{tabular}{llS[table-format=1.3]S[table-format=1.3]S[table-format=+1.3]S[table-format=2.2]}
\toprule
Backbone & Regime & {Clean depth} & {\ATK{Attacked}} & {Inflation} & {Acc.\ kept} \\
\midrule
ResNet-18   & Clean & 2.393 & 3.579 & +1.186 & 60.58 \\
VGG16       & Clean & 3.533 & 4.702 & +1.169 & 72.10 \\
MobileNetV2 & Clean & 6.556 & 7.707 & +1.151 & 97.20 \\
\midrule
ResNet-18   & Adv.  & 5.130 & 5.404 & +0.274 & 92.00 \\
VGG16       & Adv.  & 6.272 & 6.311 & +0.038 & 80.95 \\
MobileNetV2 & Adv.  & 7.208 & 7.262 & +0.054 & 67.10 \\
\bottomrule
\end{tabular}
\vspace{0.4em}
\begin{alertblock}{A consistent, architecture-independent result}
\small Clean-trained models inflate by $\approx +1.15$ steps in \emph{all three}
backbones, roughly $20\times$ the inflation of their adversarially trained counterparts,
\textbf{while the label is preserved}. Adversarial training is an unintentional defence
against energy attacks.
\end{alertblock}
\end{frame}

\begin{frame}{Training Schedule Changes Everything}\framesubtitle{Fixed-$T8$ versus mixed-$T$ training, ResNet-18}
\centering\small
\begin{tabular}{lS[table-format=+1.2]S[table-format=+1.2]}
\toprule
Regime & {Fixed-$T8$ trained} & {Mixed-$T$ trained} \\
\midrule
Clean       & +2.73 & +0.14 \\
Adversarial & +1.37 & +0.68 \\
\bottomrule
\end{tabular}
\vspace{0.6em}
\begin{block}{Why}
\small Training across $T=1 \ldots 8$ makes the \emph{short-horizon fixed-$T$ controls
much stronger} (adversarial ResNet at $T{=}4$: $69.35 \rightarrow 76.26$; at $T{=}2$:
$57.51 \rightarrow 70.04$). The VT advantage is measured against those controls, so it
collapses into the noise floor.
\end{block}
\begin{alertblock}{}
\small \textbf{Much of the reported VT benefit reflects baselines that were never
trained to run at short horizons.} The field compute-matches the baseline but never
\emph{training}-matches it.
\end{alertblock}
\end{frame}

\begin{frame}{Why Does VT Collapse Under AutoAttack?}\framesubtitle{Mechanism: hypothesis and evidence}
\begin{columns}[T]
\begin{column}{0.52\textwidth}
\textbf{Hypothesis}
\begin{enumerate}\small
  \item The exit gate is a second, \emph{low-dimensional} control channel.
  \item The attacker drives a \ATK{confident-but-wrong} early exit.
  \item Apparent PGD robustness is \ATK{gradient masking} from temporal averaging.
  \item The Square and FAB components of AutoAttack bypass this.
\end{enumerate}
\end{column}
\begin{column}{0.48\textwidth}
\textbf{Supporting evidence}
\begin{itemize}\small
  \item Deep-halting checkpoints ($\theta{=}0.85$, depth 6.2) show \emph{no} collapse.
  \item Shallow-halting ($\theta{=}0.50$, depth 4.3) shows the 23-point gap.
  \item Under PGD, accuracy falls $28.6$ points while confidence falls only $11\%$:
        the gate reads confidence, not correctness.
\end{itemize}
\end{column}
\end{columns}
\vspace{0.3em}
\centering\figorbox{confidence_trajectory.png}{0.55\linewidth}
\end{frame}

\section{Conclusions}
\begin{frame}{Limitations}\framesubtitle{Threats to validity}
\scriptsize
\begin{columns}[T]
\begin{column}{0.5\textwidth}
\begin{block}{Run-to-run variance is large}
Re-training an identical configuration moved clean accuracy by $3.51$ points and the
VT advantage from $+1.37$ to $+2.89$. Differences below roughly $2$ points are
directional, not precise.
\end{block}
\vspace{0.25em}
\begin{block}{Multi-seed bars cover one configuration}
ResNet-18 and VGG16 adversarial are each replicated over three training seeds; the
clean-trained and MobileNetV2 configurations are single-run. Attack-seed variance is measured separately (std $\le 1.33$).
\end{block}
\vspace{0.25em}
\begin{block}{Scope}
One VT policy family, one dataset (CIFAR-10).
\end{block}
\end{column}
\begin{column}{0.5\textwidth}
\begin{block}{One withheld measurement}
A full-test attack evaluation for the adversarially trained ResNet-18 could not be
reconciled with two independent subset evaluations, and is excluded pending
re-measurement.
\end{block}
\vspace{0.25em}
\begin{block}{Mixed-$T$ evidence is ResNet-only}
VGG16 and MobileNetV2 were not retrained under the mixed schedule.
\end{block}
\end{column}
\end{columns}
\end{frame}

\begin{frame}{Conclusions}\framesubtitle{What the evidence supports}
\footnotesize
\begin{columns}[T]
\begin{column}{0.5\textwidth}
\begin{block}{A real but modest advantage}
VT gives a \textbf{real but modest} clean-accuracy advantage over compute-matched
fixed-$T$: a consistent $\approx +3$ points across three backbones, \textbf{clean-trained
only}, replicated over three seeds on two backbones ($+2.13 \pm 0.85$ ResNet-18,
$+2.80 \pm 1.18$ VGG16).
\end{block}
\vspace{0.3em}
\begin{block}{It fades under harder training}
That advantage \textbf{largely disappears} under adversarial training, and under
mixed-$T$ training.
\end{block}
\end{column}
\begin{column}{0.5\textwidth}
\begin{block}{Not a robustness improvement}
VT is \textbf{not} a demonstrated robustness change in either direction. Across
training seeds the attacked comparisons are indistinguishable from zero.
\end{block}
\vspace{0.3em}
\begin{block}{A genuine attack surface}
The halting gate is a \textbf{genuine attack surface}: compute can be inflated while
accuracy is preserved, consistently across architectures.
\end{block}
\end{column}
\end{columns}
\end{frame}

\begin{frame}{Recommendations}\framesubtitle{For future adaptive-inference work}
\small
\begin{columns}[T]
\begin{column}{0.5\textwidth}
\begin{block}{Report compute-matched controls}
\FIX{Compute-matched} fixed-$T$ controls by default, never VT versus $T=8$ alone.
\end{block}
\vspace{0.3em}
\begin{block}{Training-match the baseline}
Not only compute-match it.
\end{block}
\vspace{0.3em}
\begin{block}{Use a VT-neutral criterion}
Select operating points with a \textbf{VT-neutral} criterion.
\end{block}
\end{column}
\begin{column}{0.5\textwidth}
\begin{block}{Do not stop at PGD}
Never conclude robustness from PGD alone.
\end{block}
\vspace{0.3em}
\begin{block}{Separate the two robustnesses}
Report \textbf{compute robustness} separately from label robustness.
\end{block}
\vspace{0.3em}
\begin{block}{Report seeds}
Attacked accuracy varies far more across training seeds than clean accuracy does.
\end{block}
\end{column}
\end{columns}
\end{frame}

\begin{frame}{Future Work}\framesubtitle{}
\small
\begin{columns}[T]
\begin{column}{0.5\textwidth}
\begin{block}{Multi-seed training bars}
Extend beyond VGG16 to the remaining headline configurations.
\end{block}
\vspace{0.3em}
\begin{block}{Mixed-$T$ for other backbones}
Mixed-$T$ schedules for VGG16 and MobileNetV2.
\end{block}
\vspace{0.3em}
\begin{block}{Transferable compute-overuse attacks}
Universal and transferable compute-overuse perturbations.
\end{block}
\end{column}
\begin{column}{0.5\textwidth}
\begin{block}{Timing side-channel attacks}
Black-box timing side-channel attacks on the halting decision.
\end{block}
\vspace{0.3em}
\begin{block}{Defences}
Hard compute caps, confidence smoothing, adversarial training against the compute
objective.
\end{block}
\end{column}
\end{columns}
\end{frame}
\begin{frame}{References}
\framesubtitle{References (1/3)}
\small
\setlength{\itemsep}{2pt}
\begin{thebibliography}{99}
\bibitem{graves2016}
A. Graves. Adaptive computation time for recurrent neural networks.
\emph{arXiv:1603.08983}, 2016.
\bibitem{teerapittayanon2016}
S. Teerapittayanon, B. McDanel, H. T. Kung.
BranchyNet: Fast inference via early exiting from deep neural networks.
\emph{ICPR}, 2016.
\bibitem{huang2018}
G. Huang et al.
Multi-scale dense networks for resource efficient image classification.
\emph{ICLR}, 2018.
\bibitem{li2023dynamic}
Y. Li et al.
Unleashing the potential of spiking neural networks by dynamic confidence.
\emph{ICCV}, 2023.
\bibitem{li2023input}
Y. Li et al.
Input-aware dynamic timestep spiking neural networks for efficient
in-memory computing.
\emph{DAC}, 2023.
\bibitem{chen2023}
Chen et al.
Knowing when to stop: Delay-adaptive spiking neural network classifiers
with reliability guarantees.
\emph{arXiv preprint}, 2023.
\bibitem{rao2025}
Rao et al.
VT-SNN: Variable time-step spiking neural network.
\emph{MICCAI}, 2025.
\end{thebibliography}
\end{frame}

\begin{frame}{References}
\framesubtitle{References (2/3)}
\small
\setlength{\itemsep}{2pt}
\begin{thebibliography}{99}
\bibitem{sorbaro2020}
M. Sorbaro et al.
Optimizing the energy consumption of spiking neural networks for
neuromorphic applications.
\emph{Frontiers in Neuroscience}, 14:662, 2020.
\bibitem{fang2021}
W. Fang et al.
Deep residual learning in spiking neural networks.
\emph{NeurIPS}, 2021.
\bibitem{deng2022}
S. Deng, Y. Li, S. Zhang, S. Gu.
Temporal efficient training of spiking neural network via gradient
re-weighting.
\emph{ICLR}, 2022.
\bibitem{na2022}
B. Na et al.
AutoSNN: Towards energy-efficient spiking neural networks.
\emph{ICML}, 2022.
\bibitem{sharmin2019}
S. Sharmin et al.
A comprehensive analysis on adversarial robustness of spiking neural
networks.
\emph{IJCNN}, 2019.
\bibitem{sharmin2020}
S. Sharmin et al.
Inherent adversarial robustness of deep spiking neural networks:
Effects of discrete input encoding and non-linear activations.
\emph{ECCV}, 2020.
\bibitem{liang2022}
L. Liang et al.
Toward robust spiking neural network against adversarial perturbation.
\emph{arXiv:2205.01625}, 2022.
\end{thebibliography}
\end{frame}

\begin{frame}{References}
\framesubtitle{References (3/3)}
\small
\setlength{\itemsep}{2pt}
\begin{thebibliography}{99}
\bibitem{wang2025}
J. Wang et al.
Random heterogeneous spiking neural network for adversarial defense.
\emph{iScience}, 2025.
\bibitem{ding2025}
Ding et al.
Neuromorphic computing paradigms enhance robustness through spiking
neural networks.
\emph{Nature Communications}, 2025.
\bibitem{madry2018}
A. Madry et al.
Towards deep learning models resistant to adversarial attacks.
\emph{ICLR}, 2018.
\bibitem{athalye2018}
A. Athalye, N. Carlini, D. Wagner.
Obfuscated gradients give a false sense of security:
Circumventing defenses to adversarial examples.
\emph{ICML}, 2018.
\bibitem{croce2020}
F. Croce, M. Hein.
Reliable evaluation of adversarial robustness with an ensemble of
diverse parameter-free attacks.
\emph{ICML}, 2020.
\bibitem{tramer2020}
F. Tramer, N. Carlini, W. Brendel, A. Madry.
On adaptive attacks to adversarial example defenses.
\emph{NeurIPS}, 2020.
\end{thebibliography}
\end{frame}

\appendix
\section{Backup}
\begin{frame}{Backup}\framesubtitle{}\centering\Large Backup Slides\end{frame}

\begin{frame}{B1: Selection Criterion Sensitivity}\framesubtitle{Neutral criterion at three tolerances}
\centering\small
\begin{tabular}{lS[table-format=+1.2]S[table-format=+1.2]S[table-format=+1.2]}
\toprule
Configuration & {$\delta=0.5$} & {$\delta=1.0$} & {$\delta=2.0$} \\
\midrule
R18 clean fixed-$T8$ & +1.15 & +2.73 & +2.24 \\
R18 adv.\ fixed-$T8$ & +2.03 & +1.37 & +1.37 \\
R18 clean mixed-$T$  & +0.14 & +0.14 & +0.14 \\
R18 adv.\ mixed-$T$  & +0.68 & +0.68 & -0.35 \\
\bottomrule
\end{tabular}
\\\vspace{0.4em}
\small Ordering and conclusion are stable; magnitudes are not.
\end{frame}

\begin{frame}{B2: Fixed-$T$ Sweeps by Training Schedule}\framesubtitle{ResNet-18, clean accuracy (\%)}
\centering\scriptsize
\begin{tabular}{lcccccccc}
\toprule
Configuration & $T{=}1$ & $T{=}2$ & $T{=}3$ & $T{=}4$ & $T{=}5$ & $T{=}6$ & $T{=}7$ & $T{=}8$ \\
\midrule
clean fixed-$T8$ & 10.00 & 84.72 & 87.89 & 90.11 & 90.48 & 90.97 & 91.09 & 91.26 \\
clean mixed-$T$  & 41.39 & 81.48 & 86.75 & 88.78 & 89.80 & 90.26 & 90.61 & 90.57 \\
adv.\ fixed-$T8$ & 12.11 & 57.51 & 66.79 & 69.35 & 71.46 & 72.98 & 73.23 & 73.67 \\
adv.\ mixed-$T$  & 46.28 & 70.04 & 74.31 & 76.26 & 76.82 & 77.14 & 77.58 & 77.79 \\
\bottomrule
\end{tabular}
\\\vspace{0.4em}
\small Mixed-$T$ lifts adversarial training at \emph{every} horizon; for clean training
it only rescues $T=1$.
\end{frame}

\begin{frame}{B3: Attack-Seed Variance}\framesubtitle{Three attack seeds, one checkpoint, $n=1000$}
\centering\small
\begin{tabular}{llc}
\toprule
Operating point & Attack & Mean $\pm$ std \\
\midrule
\FIX{Fixed $T=8$}   & PGD        & $48.83 \pm 0.31$ \\
\FIX{Matched fixed} & PGD        & $48.07 \pm 0.85$ \\
\VT{VT selected}    & PGD        & $49.03 \pm 0.50$ \\
\midrule
\FIX{Fixed $T=8$}   & AutoAttack & $31.53 \pm 0.75$ \\
\FIX{Matched fixed} & AutoAttack & $30.00 \pm 0.85$ \\
\VT{VT selected}    & AutoAttack & $29.83 \pm 1.33$ \\
\bottomrule
\end{tabular}
\\\vspace{0.4em}
\small Attack randomness contributes at most $1.33$ points. Training variance is larger.
\end{frame}

\begin{frame}{B4: Clean-Trained Models Collapse Entirely}\framesubtitle{No robustness to preserve}
\centering\small
\begin{tabular}{lS[table-format=2.2]S[table-format=2.2]S[table-format=2.2]S[table-format=2.2]}
\toprule
Backbone & {PGD \FIX{fixed}} & {PGD \VT{VT}} & {AA \FIX{fixed}} & {AA \VT{VT}} \\
\midrule
ResNet-18   & 0.00 & 0.02 & 0.06 & 0.02 \\
VGG16       & 0.05 & 0.10 & 0.00 & 0.00 \\
MobileNetV2 & 5.10 & 4.95 & 5.25 & 5.40 \\
\bottomrule
\end{tabular}
\\\vspace{0.4em}
\small The inference policy is irrelevant when the backbone has no adversarial training.
\end{frame}

\begin{frame}{B5: SynOps Accounting}\framesubtitle{Methodology}
\begin{itemize}\small
  \item SynOps $=$ dense ops $\times$ input activation sparsity, summed over convolutional
        and linear layers, over all executed timesteps.
  \item \texttt{avg\_dense\_ops} depends only on architecture and $T$, and is used here as
        an architecture fingerprint.
  \item VT compute is measured \textbf{sample-exact}: each sample is executed individually,
        not as a batched proxy.
  \item SynOps is an algorithmic proxy for energy, \textbf{not} a hardware measurement.
\end{itemize}
\end{frame}

\end{document}
